\subsection{Collection before the soft endpoint evaluation}
The soft coefficient can be shortened further by collecting the rational
integration-by-parts terms before expanding the logarithms into GPLs.
This also removes separate divergent boundary terms. Keep $X=s-1>0$,
$\beta=X/(X+2)$, and use the original angular variable $0<U<1$.
Define
\begin{align}
 k&=\sqrt{\frac{1+XU^2}{s}},\qquad d=1-\beta^2U^2,\\
 a&=\operatorname{artanh}(\beta U),\quad
 L=\operatorname{artanh}U,\quad f=\operatorname{artanh}(U/k).
\end{align}
The square root is positive. Since $k^2-U^2=(1-U^2)/s>0$, both
inverse hyperbolic tangents involving $U$ are real in the open interval.
Their derivatives are
\begin{equation}
 a'=\frac{\beta}{d},\qquad L'=\frac1{1-U^2},\qquad
 f'=\frac1{k(1-U^2)},\qquad k'=\frac{XU}{sk}.
\end{equation}
Primes here mean differentiation with respect to $U$.

\paragraph{Rational integration by parts.} Let $\Phi(U)$ be the soft
integrand after the pole removal in
\eqref{eq:comp-ck-pole-removal}, including its common factor
$4/(\beta E^2)$. Rational division and the derivatives above give
\begin{equation}
 \Phi(U)=\frac{\dd\mathcal B}{\dd U}+\sum_i T_i(U)f_i(U),
 \qquad (f_i)=(aL,af,a,L,f,1),
\end{equation}
where the rational boundary can be chosen as
\begin{equation}
 \mathcal B=Q_{aL}aL+Q_{af}af+Q_a a+Q_L L+Q_1.
\end{equation}
Its coefficients are
\begin{align}
 Q_{aL}&=\frac{2U(X+2)}{X^3s}
       \bigl(U^2X^3+X^3+24X^2+36X+16\bigr),\\
 Q_{af}&=\frac{128s^2Uk}{X^3(X+4)d},\\
 Q_a&=-\frac{2U^2(X+2)}s,\qquad
 Q_L=\frac{U^2(X+2)^2}{Xs},\\
 Q_1&=-\frac{U(X+2)^2}{Xs}.
\end{align}
For example, the root-dependent part follows directly from
\begin{equation}
 \frac{\dd}{\dd U}\frac{Uk}{d}
 =\frac{1+(2X+\beta^2)U^2}{skd^2}.
\end{equation}
Matching its constant and $U^2$ coefficients removes the double pole
from the coefficient of $af$. Substituting $f'$ then gives a rational
coefficient of $a$, because the factor $k$ in $Q_{af}$ cancels.
The other descendants follow by differentiating $aL$, $a$, and $L$.
Each step is a rational identity after replacing $k^2$ by
$(1+XU^2)/s$.

\paragraph{Vanishing logarithmic boundaries.} Replace each of
$Q_{aL},Q_{af},Q_a,Q_L$ by $Q_i(U)-Q_i(1)$ and add
$Q_i(1)f_i'$ to the remainder. This leaves the full integrand unchanged.
All these differences are $O(1-U)$. At the upper endpoint $a$ is finite,
while $L$ and $f$ grow logarithmically. Every logarithmic boundary
therefore vanishes. At $U=0$, all three logarithms vanish. The only
integrated boundary is
\begin{equation}
 Q_1(1)-Q_1(0)=-\frac{(X+2)^2}{Xs}.
\end{equation}
The cancellation of the apparent pole in the coefficient of $a$ uses
\begin{equation}
 \frac{1/k-1}{1-U^2}=\frac{X}{sk(1+k)}.
\end{equation}
This identity follows by multiplying by $1+k$ and using the displayed
expression for $k^2$. It also fixes the finite endpoint value without
subtracting two separately divergent numerical quantities.

Define the seven energy-dependent coefficients
\begin{align}
 A_L&=\frac{4(X+2)(X^3+12X^2+18X+8)}{X^3s},\\
 A_f&=\frac{32s(X+2)^2}{X^3(X+4)},\\
 P&=-\frac{16s(X^2-4X-6)}{X^2(X+2)},\\
 B&=\frac{16s(X^4+2X^3-24X^2-80X-48)}{X^3(X+2)(X+4)},\\
 C_L&=-\frac{4(X^3+14X^2+20X+8)}{X^2s},\\
 M&=-\frac{32s(X+2)}{X^2(X+4)},\qquad
 C_0=\frac{8(X+2)}{X^2}.
\end{align}
Here $A_L=Q_{aL}(1)$ and $A_f=Q_{af}(1)$.
The resulting remainder is the short expression
\begin{equation}
\begin{aligned}
 \widehat\Phi(U)={}&\frac1d\biggl[
 PaL+\frac{Baf}{k}
 +(A_L+A_f)\frac{1+\beta^2U}{1+U}a\\
 &\hspace{12mm}+(C_LU+A_L\beta)L
 +\left(\frac{MU}{k}+A_f\beta\right)f+C_0\beta\biggr]
 +\frac{A_fXa}{sk(1+k)}.
\end{aligned}
\label{eq:comp-ck-regular-root-kernel}
\end{equation}
None of its rational coefficients has a pole on the closed physical
interval. The logarithmic growth at $U=1$ is integrable. In particular,
the constant-function term integrates immediately to
\begin{equation}
 \int_0^1\frac{C_0\beta}{d}\,\dd U
 =C_0\operatorname{artanh}\beta
 =\frac{4(X+2)}{X^2}\ln s.
\end{equation}
Thus the already integrated soft coefficient admits the exact
reorganization
\begin{equation}
 \mathcal C_K(s)=-\frac{(X+2)^2}{Xs}
 +\int_0^1\widehat\Phi(U)\,\dd U
 -\frac{\operatorname{artanh}\beta}{\beta}\mathcal C_0(s).
\label{eq:comp-ck-regular-root-identity}
\end{equation}
Equation~\eqref{eq:comp-ck-regular-root-identity} is an identity for
collecting the existing analytic endpoint expression. Its remaining
terms are evaluated with the logarithmic primitives of the preceding
subsection. Under $U=(w-w^{-1})/(2\sqrt X)$, all their coefficients
are proper rational functions with simple poles, so only the single-pole
primitives are needed. The square root has not been expanded into
additional energy-dependent letters at the collection step.

\paragraph{Finite analytic endpoint sum.} The reorganized coefficient
can also be stated without a remaining integral. Put
\begin{equation}
 W=\sqrt s+\sqrt X,\qquad Z=W-1,\qquad
 r=\sqrt{1+X/\beta^2}-\frac{\sqrt X}{\beta}.
\end{equation}
Define three lists of pairs consisting of a letter and its weight,
\begin{align}
 \mathcal A&=\{(r,\tfrac12),(-r^{-1},\tfrac12),
                  (r^{-1},-\tfrac12),(-r,-\tfrac12)\},\\
 \mathcal L&=\{(W^{-1},\tfrac12),(-W,\tfrac12),
                  (W,-\tfrac12),(-W^{-1},-\tfrac12)\},\\
 \mathcal F&=\{(W^{-1},\tfrac12),(-W^{-1},\tfrac12),
                  (W,-\tfrac12),(-W,-\tfrac12)\}.
\end{align}
For example, $a(w)=\sum_{(b,h)\in\mathcal A}
h\ln[(w-b)/(1-b)]$. The other two lists give $L(w)$ and $f(w)$.
Use $U=(w-w^{-1})/(2\sqrt X)$ and $k=(w+w^{-1})/(2\sqrt s)$
in the five rational functions
\begin{align}
 R_{aL}&=U'\frac P d,\qquad R_{af}=U'\frac B{kd},\\
 R_a&=U'\left[\frac{(A_L+A_f)(1+\beta^2U)}{(1+U)d}
                      +\frac{A_fX}{sk(1+k)}\right],\\
 R_L&=U'\frac{C_LU+A_L\beta}{d},\qquad
 R_f=U'\frac{MU/k+A_f\beta}{d}.
\end{align}
Here $U'=\dd U/\dd w$. Each function is proper rational and has
only simple poles. If $\mathcal P_R$ is its pole set, its coefficients
are given by the finite algebraic limits
\begin{equation}
 R(w)=\sum_{c\in\mathcal P_R}\frac{\rho_c}{w-c},\qquad
 \rho_c=\lim_{w\to c}(w-c)R(w).
\end{equation}
For a list $\mathcal B$ and two lists $\mathcal B,\mathcal D$, define
the following finite sums of ordinary GPLs:
\begin{align}
 \mathcal S[R;\mathcal B]
 &=\sum_{c\in\mathcal P_R}\sum_{(b,h)\in\mathcal B}
       \rho_ch\,G(c-1,b-1;Z),\\
 \mathcal S[R;\mathcal B,\mathcal D]
 &=\sum_{c\in\mathcal P_R}
   \sum_{(b,h)\in\mathcal B}\sum_{(e,j)\in\mathcal D}
   \rho_chj\bigl[G(c-1,b-1,e-1;Z)
                    +G(c-1,e-1,b-1;Z)\bigr].
\end{align}
These follow by substituting the partial fractions and using the shuffle
identity for the product of the two logarithms. The soft coefficient is
therefore
\begin{equation}
\begin{aligned}
 \mathcal C_K(s)={}&-\frac{(X+2)^2}{Xs}
 +\frac{4(X+2)}{X^2}\ln s
 -\frac{\ln s}{2\beta}\mathcal C_0(s)\\
 &+\mathcal S[R_{aL};\mathcal A,\mathcal L]
  +\mathcal S[R_{af};\mathcal A,\mathcal F]\\
 &+\mathcal S[R_a;\mathcal A]
  +\mathcal S[R_L;\mathcal L]
  +\mathcal S[R_f;\mathcal F].
\end{aligned}
\label{eq:comp-ck-collected-finite-gpl}
\end{equation}
All the residues in this expression are algebraic functions specified
by the five short rational functions above. No endpoint extrapolation
or numerical integration is part of its definition.
The possible poles belong to
$\{0,\pm r,\pm r^{-1},W^{-1},-W,-W^{-1}\}$.
Since $0<r<W^{-1}<1<W<r^{-1}$, no outer letter lies on the
integration interval. The inner letter $W$ produces only the
integrable endpoint logarithms. Thus every GPL in this endpoint sum
is finite separately, and all logarithms follow the real physical
branch established above.
